\documentclass[12pt]{article}
\usepackage[utf8]{inputenc}
\usepackage[T1]{fontenc}
\usepackage{hyperref}
\title{Viemo Language Resources for K-12 Curriculum Development}
\author{Compiled by Claude Sonnet 5 (claude-sonnet-5) for Andr\'as Kornai}
\date{}
\begin{document}
\maketitle

\section{Introduction}
Viemo (ISO 639-3 \texttt{vig}; also Vige, Vigu\'e, Vigye) is spoken in southwestern Burkina Faso, southeast of Bobo-Dioulasso. Wikipedia cites about 8,000 speakers (1995); a Global Recordings page gives 12,500. Its exact Niger-Congo affiliation is unresolved, and the community is highly bilingual in Jula. Very little educational material was found in search results.

\section{Primary Educational Resources}
No primers, textbooks or school programmes for Viemo were found in the search results.

\section{Linguistic and Reference Materials}
\begin{itemize}
\item Viemo language (Wikipedia overview): \url{https://en.wikipedia.org/wiki/Viemo_language}
\item Other Languages page, dogonlanguages.info (listing surfaced in search; contents not verified): \url{https://dogonlanguages.info/other}
\end{itemize}

\section{Religious and Community Resources}
\begin{itemize}
\item Joshua Project, Viemo in Burkina Faso: \url{https://joshuaproject.net/people_groups/15790}
\item Scripture Earth listing for vig (contents not verified): \url{https://www.scriptureearth.org/00eng.php?iso=vig}
\end{itemize}

\section{Digital and Audio Resources}
\begin{itemize}
\item Global Recordings Network, language page for vig: \url{https://globalrecordings.net/en/language/vig}
\end{itemize}

\section{Curriculum Development Context}
Viemo is not known to be taught in schools, and the search found no orthography, grammar or dictionary. A curriculum would need new documentation work, probably with Jula and French as bridging languages.
\end{document}